% Generated draft: requires mathematical and source review.
\documentclass[11pt]{article}
\usepackage[margin=25mm]{geometry}
\usepackage{fontspec}
\setmainfont{Latin Modern Roman}
\usepackage{amsmath,amssymb,mathtools,unicode-math}
\setmathfont{Latin Modern Math}
\usepackage{enumitem,needspace,xurl,hyperref}
\setlength{\parindent}{0pt}
\setlength{\parskip}{5pt}
\setlength{\emergencystretch}{3em}
\newcommand{\N}{\mathbb{N}}
\newcommand{\Z}{\mathbb{Z}}
\newcommand{\Q}{\mathbb{Q}}
\newcommand{\R}{\mathbb{R}}
\newcommand{\C}{\mathbb{C}}
\newcommand{\F}{\mathbb{F}}
\newcommand{\A}{\mathbb{A}}
\newcommand{\PP}{\mathbb{P}}
\newcommand{\E}{\mathbb{E}}
\newcommand{\eps}{\varepsilon}
\newcommand{\dd}{\,\mathrm{d}}
\begin{document}
\section*{UnsolvedMath ID: 11100057}
\textbf{Source problem code:} catalogue entry
\subsection{Definitions and mathematical statement}
The source asks to compare Rezk's homotopy theory of homotopy theories with Hovey's proposed $2$-category of model categories, and to make precise an equivalence between the two viewpoints. Neither the morphisms and $2$-morphisms in the proposed $2$-category nor the meaning of equivalence is supplied in this item.

For the Rezk viewpoint named by the later summary, a simplicial space is a functor $W:\Delta^{\mathrm{op}}\to\mathrm{sSet}$. A complete Segal space is a Reedy fibrant such $W$ whose Segal maps
\[
 W_n\longrightarrow W_1\mathbin{\times}_{W_0}\cdots\mathbin{\times}_{W_0}W_1\qquad(n\geq2)
\]
are weak equivalences and whose degeneracy map $W_0\to W_{\mathrm{heq}}$ is a weak equivalence, where $W_{\mathrm{heq}}\subseteq W_1$ consists of components representing equivalences. These conditions encode composition associative up to coherent homotopy rather than a strictly specified composition on points.

For a model category $\mathcal M$ with weak equivalences $\mathcal W$, its Dwyer--Kan simplicial localization $L(\mathcal M,\mathcal W)$ is a candidate presentation of its underlying homotopy theory. The requested comparison should specify which model categories and morphisms are admitted, which equivalences are inverted, the corresponding Rezk objects, and the categorical level at which a comparison is an equivalence. Size universes or another treatment of large categories are also required. The source imposes no accessibility, combinatorial, or presentability assumptions, and does not specify a restricted class of Rezk objects; these cannot be silently added. In particular, the later summary's broad equivalence claim about a homotopy/$2$-category of model categories does not by itself supply a precise comparison theorem.
\subsection{Short English statement}
Explain in a precise sense how Rezk's homotopy theories and a suitable $2$-category of model categories describe the same information. The supplied summary reports modern equivalences between models for higher categories and labels the problem solved. The precise model-category side, its size restrictions, and the level of equivalence remain absent from this isolated source.
\subsection{Sources}
\begin{itemize}
\item Local UnsolvedMath registry, permanent ID $11100057$, source code \texttt{catalogue entry}, proposed by Mark Hovey; created 2026-07-29 and updated 2026-08-06. Supplied status \texttt{solved}; classification \texttt{SOLVED-IN-LITERATURE}.
\item Mark Hovey, \emph{Hovey's Algebraic Topology Problem List}, ``Model categories,'' item 11, surviving mirror \url{https://www-users.cse.umn.edu/~tlawson/hovey/}, accessed 2026-07-29; HTML ordered-list extraction. Background: 65 top-level items in nine category pages; default difficulty L3. The original names Charles Rezk, emphasizes the largeness of the category, and proposes equivalence ``in some sense.''
\item The research summary attributes complete Segal spaces to Rezk and cites Joyal--Tierney, Bergner, Lurie, and Dwyer--Kan in discussing complete Segal spaces, quasicategories, simplicial categories, simplicial localization, and the $\infty$-category of $\infty$-categories. It asserts a comparison with a homotopy/$2$-category of model categories, but gives no exact references or hypotheses. Its final discussion of what Dwyer--Kan localization sends a model category to is truncated.
\item Background status evidence \texttt{NEEDS\_REVIEW}: the source presents an open problem and a current primary-literature audit remains. Separate rights note \texttt{NEEDS\_REVIEW}: publicly accessible author/publisher copy with redistribution terms requiring review. Dataset attribution: Ulam AI / UnsolvedMath Contributors, CC BY 4.0.
\end{itemize}
\end{document}
